\section{Conclusion}
\label{sec:conclusion}

We began with the observation that most machine learning datasets are invisible --- not because they lack quality, but because they lack tags. This paper has demonstrated that the invisibility is quantifiable, structural, and remediable.

\subsection{Summary of Contributions}

\begin{enumerate}
\item \textbf{Binary tag presence yields a 22.6\% adoption advantage} (IRR=1.2263, 95\% CI [1.1681, 1.2873], $p<0.001$, N=5,217) that is structural --- surviving decade fixed effects with only 12.2\% attenuation despite extreme era-tag collinearity (Cram\'er's V=0.823).

\item \textbf{Tag count magnitude amplifies the advantage log-linearly} within tagged datasets (IRR=1.5332 per $\log(\text{tag\_count}+1)$, N=2,625, $p<0.001$), with essentially zero decade confounding (attenuation ratio=1.0007).

\item \textbf{Comprehensive tagging (6+ tags) is the dominant categorical tier} (IRR=1.2861, meaningfully distinct from 1--5 tags at $p=5.54 \times 10^{-10}$), while OpenML tagging behavior is bimodal --- users tag comprehensively or not at all.
\end{enumerate}

The informative negative result (H-M3) honestly documents the limits of the categorical characterization while refining the practical recommendation toward the 6+ tag target.

\subsection{Future Directions}

\textbf{From untested alternatives:} Reverse causality testing via OpenML API timestamp data --- whether popular datasets attract tags retroactively --- requires tag assignment timestamps unavailable in the current corpus. A natural experiment around API version changes affecting tagging permissions would provide stronger causal evidence.

\textbf{From unverified assumptions:} OpenML search API log analysis (query $\rightarrow$ click $\rightarrow$ task creation sequences) would directly verify the tag-indexed search pathway mechanism, converting proxy evidence (IRR dose-response consistent with mechanism) to direct evidence.

\textbf{From scope extensions:} Multi-platform replication (HuggingFace, Kaggle, UCI) would test whether the FAIR F1 threshold-plus-amplification structure generalizes across repository architectures. A hurdle model on the full corpus (including $N\_tasks=0$) would estimate $P(\text{any adoption})$ alongside conditional adoption intensity.

\subsection{Closing}

As ML repositories grow and dataset proliferation accelerates, the structural advantages of keyword-indexed FAIR F1 compliance will compound. The ML community already struggles to find datasets appropriate for specific tasks; this challenge will not diminish as catalogs grow. A simple act at dataset upload --- adding 6 or more descriptive keyword tags --- is associated with substantially higher conditional adoption. While causal ordering requires timestamp verification beyond the available corpus, the magnitude of the association and its structural robustness across model variants make keyword tagging a well-motivated metadata priority. We hope this quantification motivates both individual creators and repository designers to treat FAIR F1 tagging not as a documentation checkbox, but as a discoverability mechanism with measurable adoption consequences.
